\subsection{Finite conjugation-fixed class modules and source-group freeness}
\label{subsec:cm-finite-plus}

Continue with \(K,p,L_n,L_\infty,\Omega,H,A\) from
Subsection~\ref{subsec:cm-elliptic-radical}.
Complex conjugation is denoted by \(c\).
Put
\[
 \mathcal X_2=\varprojlim_n\operatorname{Cl}(L_n)[p^\infty],
 \qquad
 \mathcal J_2=\varinjlim_n\operatorname{Cl}(L_n)[p^\infty],
\]
where the first transition maps are ideal norms and the second
are extension of ideals.
The finite-level class groups here are the full class groups.
Let \(\mathcal X_2^+=(1+c)\mathcal X_2/2\), considered as a compact
\(\mathbb Z_p\)-module.
This definition does not assume that the fixed subgroup is stable
under the two-variable Iwasawa algebra.
Write
\[
 G=\operatorname{Gal}(L_\infty/K),\quad
 \Gamma=\operatorname{Gal}(L_\infty/L_1)\cong\mathbb Z_p^2,\quad
 L_{\mathrm c}=L_1(\mu_{p^\infty}),\quad
 \Gamma_{\mathrm a}=\operatorname{Gal}(L_\infty/L_{\mathrm c})
                   \cong\mathbb Z_p .
\]
Ray reciprocity identifies
\[
 G=\bigl(\mathbb Z_p^\times\times\mathbb Z_p^\times\bigr)/O_K^\times,
 \qquad
 c(x,y)c^{-1}=(y,x),\qquad
 \chi_{\mathrm{cyc}}(x,y)=xy.
\]
The embedded \(O_K^\times\) has the form
\((\varepsilon,\varepsilon^{-1})\).
Thus \(G/\langle g(g^c)^{-1}:g\in G\rangle\) is precisely the
cyclotomic quotient \(\mathbb Z_p^\times\), and
\(\Gamma_{\mathrm a}\) is its pro-\(p\) kernel.

\begin{lemma}\label{lem:cm-full-class-kummer}
Let
\[
 \mathcal U_\infty=O_{L_\infty}^{\times}
       =\bigcup_n O_{L_n}^{\times},\qquad
 L_{\mathrm u}=L_\infty(u^{1/p^r}:
                         u\in\mathcal U_\infty,\ r\geq1).
\]
There is a \(G\)- and \(c\)-equivariant exact sequence
\begin{equation}\label{eq:cm-full-class-kummer}
 0\longrightarrow
 \mathcal U_\infty\otimes_{\mathbb Z}\mathbb Q_p/\mathbb Z_p
 \longrightarrow
 B:=\operatorname{Hom}_{\mathrm{cont}}(A,\mu_{p^\infty})
 \xrightarrow{\vartheta}\mathcal J_2
 \longrightarrow0 .
\end{equation}
Consequently the restriction map has kernel
\begin{equation}\label{eq:cm-unit-annihilator}
 N_{\mathrm u}:=
 \ker\!\left(A\longrightarrow
            \operatorname{Gal}(L_{\mathrm u}/L_\infty)\right)
 \cong\operatorname{Hom}(\mathcal J_2,\mu_{p^\infty}).
\end{equation}
\end{lemma}

\begin{proof}
We spell out \(\vartheta\); this also proves that the right-hand term
is the full direct-limit class group.
The same construction for a cyclotomic tower is
\cite[Proposition 3.3.9 and Corollary 3.3.14]{SharifiIwasawaNotes}.
Its valuation proof applies to the present tower because it
contains the cyclotomic tower.

Represent a finite-order character in \(B\) by a Kummer class
\(u\in L_\infty^\times/L_\infty^{\times p^r}\).
At every prime outside \(p\), its valuation is divisible by \(p^r\),
because the corresponding Kummer extension is unramified there.
Pass to a finite ray layer containing \(u\) and sufficiently high
in the cyclotomic direction that all valuations at primes over
\(p\) are divisible by \(p^r\) as well.
This is possible since their ramification indices grow without
bound.
At that layer there is a unique fractional ideal \(\mathfrak b\)
with
\[
                         (u)=\mathfrak b^{p^r}.
\]
Set \(\vartheta(u)=[\mathfrak b]\in\mathcal J_2\).
Replacing \(u\) by \(uv^{p^r}\) changes \(\mathfrak b\) by
the principal ideal \((v)\); going to a higher layer extends
\(\mathfrak b\).
Representing the same character at exponent \(p^{r+1}\) replaces
\(u\) by \(u^p\) and leaves \(\mathfrak b\) unchanged.
Thus the map is well defined and equivariant.

Every \(p^r\)-torsion ideal class at a finite layer has a
representative \(\mathfrak b\) with
\(\mathfrak b^{p^r}=(u)\); the associated Kummer extension is
unramified outside \(p\).
This proves surjectivity.
The class \([\mathfrak b]\) dies in the direct limit exactly when,
at a higher finite layer, \(\mathfrak b=(v)\).
Then \(uv^{-p^r}\) is a global unit.
This proves the asserted kernel.
The left Kummer map is injective because a \(p^r\)-th root of a
global unit, whenever it belongs to \(L_\infty\), is again a global
unit.
Taking the Pontryagin dual with values in \(\mu_{p^\infty}\)
gives \eqref{eq:cm-unit-annihilator}.
The Tate twist in this duality is retained: \(c\) acts by \(-1\)
on \(\mu_{p^\infty}\).
\end{proof}

\begin{lemma}\label{lem:cm-anticyclotomic-invariants}
Without a class-group finiteness assumption, one has
\[
 A^{\Gamma_{\mathrm a}}=0,\qquad
 H^2(H,\mathbb Q_p/\mathbb Z_p)=0 .
\]
\end{lemma}

\begin{proof}
Put \(H_{\mathrm c}=\operatorname{Gal}(\Omega/L_{\mathrm c})\).
These are the full maximal pro-\(p\) Galois groups unramified outside
\(p\) of their indicated infinite base fields.
To check the comparison, a finite Galois \(p\)-extension over either
infinite field descends to a finite subfield of that base,
which is a \(p\)-extension of \(L_1\).
After increasing that finite subfield, take the descended
Galois \(p\)-extension and adjoin all its conjugates over \(L_1\);
the resulting normal closure has \(p\)-power degree and is
unramified outside \(p\), and therefore belongs to \(\Omega\).
Thus we have the exact sequence of arithmetic groups
\begin{equation}\label{eq:cm-anticyclotomic-group-extension}
 1\longrightarrow H\longrightarrow H_{\mathrm c}
                 \longrightarrow\Gamma_{\mathrm a}
                 \longrightarrow1 .
\end{equation}

Iwasawa's theorem on weak Leopoldt for a cyclotomic
\(\mathbb Z_p\)-extension gives
\[
 H^2(H_{\mathrm c},\mathbb Q_p/\mathbb Z_p)=0;
\]
see \cite[Remark 4.3, author version p.~26]{VenjakobIwasawa},
which explicitly states the result for every number field and
refers to its cohomological formulation.
This theorem concerns the trivial coefficient
\(\mathbb Q_p/\mathbb Z_p\).
If it is formulated first for the full \(S\)-ramified Galois
group, its inflation to the maximal pro-\(p\) quotient in degree
two is injective.
Indeed, an invariant nonzero homomorphism from the kernel of
the maximal pro-\(p\) quotient to \(\mathbb F_p\) would produce
a larger pro-\(p\) quotient; the same assertion for
\(\mathbb Q_p/\mathbb Z_p\) follows by taking a nonzero
order-\(p\) multiple.
The five-term sequence proves the required injectivity.

Let \(M=H^1(H,\mathbb Q_p/\mathbb Z_p)=A^\vee\).
Since \(\operatorname{cd}_p(\Gamma_{\mathrm a})=1\),
Hochschild--Serre for
\eqref{eq:cm-anticyclotomic-group-extension} gives
\begin{equation}\label{eq:cm-anticyclotomic-HS}
 0\longrightarrow H^1(\Gamma_{\mathrm a},M)
 \longrightarrow H^2(H_{\mathrm c},\mathbb Q_p/\mathbb Z_p)
 \longrightarrow
 H^2(H,\mathbb Q_p/\mathbb Z_p)^{\Gamma_{\mathrm a}}
 \longrightarrow0 .
\end{equation}
For a generator \(\gamma\) of \(\Gamma_{\mathrm a}\),
\[
 H^1(\Gamma_{\mathrm a},M)=M/(\gamma-1)M,
 \qquad
 \bigl(M/(\gamma-1)M\bigr)^\vee=A^{\Gamma_{\mathrm a}}.
\]
The first vanishing follows.
Equation \eqref{eq:cm-anticyclotomic-HS} also makes the
\(\Gamma_{\mathrm a}\)-invariants of the indicated \(H^2(H,-)\)
zero.
A nonzero discrete \(p\)-primary module under a pro-\(p\) group
always has nonzero invariants: choose a nonzero element of order
\(p\), take the finite-dimensional span of its finite orbit,
and use the nonzero fixed subspace of a finite \(p\)-group
acting on a nonzero \(\mathbb F_p\)-space.
Thus \(H^2(H,\mathbb Q_p/\mathbb Z_p)=0\) as well.
No freeness of \(H\) was used.
\end{proof}

\begin{proposition}\label{prop:cm-finite-plus-free}
For the preceding CM ray tower, suppose that the explicitly
defined group \(\mathcal X_2^+\) is finite.
Then
\[
 \mathcal J_2=0,\qquad
 \operatorname{Gal}(L_{\mathrm u}/L_\infty)=A,
 \qquad H\text{ is a free pro-}p\text{ group}.
\]
In particular every finite-layer \(p\)-primary ideal class becomes principal
at some higher ray layer, and the maximal abelian subextension
of \(\Omega/L_\infty\) is generated by \(p\)-power roots of
global units.
\end{proposition}

\begin{proof}
First, every norm
\[
 \operatorname{Cl}(L_{n+1})[p^\infty]
           \longrightarrow\operatorname{Cl}(L_n)[p^\infty]
\]
is surjective.
Ray reciprocity identifies
\(\operatorname{Gal}(L_{n+1}/L_n)\) with the product of the
two principal-unit quotients, each of order \(p\).
The inertia at \(\mathfrak p\) contains the first factor and
the inertia at \(\bar{\mathfrak p}\) contains the second.
Hence \(L_{n+1}/L_n\) has no nontrivial everywhere-unramified
subextension.
By the Artin reciprocity description of the image of the ideal
norm, this is exactly the required surjectivity on class groups.
The same holds for all composite norm maps.
All maps commute with \(c\); averaging by \((1+c)/2\) proves
surjectivity on their fixed subgroups.

It follows that \(\mathcal X_2^+\) surjects onto every
\(\operatorname{Cl}(L_n)[p^\infty]^+\).
Since \(\mathcal X_2^+\) is finite, their orders are eventually
constant, and for all sufficiently large \(m\geq n\) the norm
\[
 N_{m/n}:\operatorname{Cl}(L_m)[p^\infty]^+
                       \xrightarrow{\sim}
                         \operatorname{Cl}(L_n)[p^\infty]^+
\]
is an isomorphism.
Let \(j_{m/n}\) be extension of ideals.
The exact identity
\[
                  N_{m/n}j_{m/n}=[L_m:L_n]
                               =p^{2(m-n)}
\]
holds on \(\operatorname{Cl}(L_n)\).
For fixed sufficiently large \(n\), this multiplication kills
the finite group \(\operatorname{Cl}(L_n)[p^\infty]^+\) when
\(m\) is sufficiently large.
Injectivity of \(N_{m/n}\) on the plus subgroup then gives
\(j_{m/n}=0\) there.
Consequently
\begin{equation}\label{eq:cm-direct-limit-plus-zero}
                         \mathcal J_2^+=0.
\end{equation}
This argument uses the exact norm and extension maps and does not
regard a plus subgroup as a two-variable Iwasawa submodule.

By \eqref{eq:cm-unit-annihilator} and the Tate twist in that
pairing, \eqref{eq:cm-direct-limit-plus-zero} says that \(c\)
acts as \(+1\) on \(N_{\mathrm u}\).
The subgroup \(N_{\mathrm u}\) is stable under \(G\), because
global units and their radical field are stable under \(G\).
For \(g\in G\) and \(a\in N_{\mathrm u}\), one therefore has
\[
                    ga=c(ga)=g^c\,ca=g^c a.
\]
Hence every \(g(g^c)^{-1}\) acts trivially on \(N_{\mathrm u}\).
The explicit ray-reciprocity quotient preceding the lemmas shows
that \(N_{\mathrm u}\subseteq A^{\Gamma_{\mathrm a}}\).
Lemma~\ref{lem:cm-anticyclotomic-invariants} makes this group zero.
Thus \(N_{\mathrm u}=0\).
The perfect Kummer duality
\eqref{eq:cm-unit-annihilator} now also yields
\(\mathcal J_2=0\), and proves the assertion about the maximal
abelian subextension.

It remains to prove freeness.
Equation \eqref{eq:cm-full-class-kummer} and
\(\mathcal J_2=0\) identify
\[
 H^1(H,\mu_{p^\infty})
       =\mathcal U_\infty\otimes_{\mathbb Z}
                                     \mathbb Q_p/\mathbb Z_p.
\]
The group on the right is \(p\)-divisible.
Since \(H\) acts trivially on \(\mu_{p^\infty}\), choosing
compatible roots of unity identifies it with
\(H^1(H,\mathbb Q_p/\mathbb Z_p)\).
The cohomology sequence for
\[
 0\longrightarrow\mathbb F_p
  \longrightarrow\mathbb Q_p/\mathbb Z_p
  \xrightarrow{p}\mathbb Q_p/\mathbb Z_p
  \longrightarrow0
\]
therefore gives \(H^2(H,\mathbb F_p)=0\), using the already
proved vanishing of \(H^2(H,\mathbb Q_p/\mathbb Z_p)\).
The cohomological freeness criterion for pro-\(p\) groups proves
that \(H\) is free.
For the same criterion in this arithmetic setting, see
\cite[Theorem 4.17 and its proof, author version p.~33]
{VenjakobIwasawa}.
\end{proof}

The group whose freeness has just been proved is the
source \(H=\operatorname{Gal}(\Omega/L_\infty)\).
To obtain a statement about an outer quotient, one must still
apply the specific restriction map defined by \(\rho\).
The radical calculation of the preceding subsection supplies
that map on the abelian quotient once the finite-level
elliptic-unit/global-unit index is compared.

\subsubsection{Character components and the exact Kummer twist}

Let
\[
 \Delta=\operatorname{Gal}(L_1/K)
       \cong\bigl((\mathbb F_p^\times)^2\bigr)/O_K^\times .
\]
We regard \(\Delta\) as the canonical prime-to-\(p\) factor of
\(G=\Delta\times\Gamma\).
All its characters take values in \(\mathbb Z_p^\times\), since
its exponent divides \(p-1\).
For such a character \(\chi\), write
\[
 e_\chi=\frac1{|\Delta|}
             \sum_{\delta\in\Delta}\chi(\delta)^{-1}\delta,
 \qquad
 \chi^c(\delta)=\chi(c\delta c^{-1}),\qquad
 \chi^\dagger=\omega_{\mathrm{cyc}}\chi^{-1},
\]
where \(\omega_{\mathrm{cyc}}=\chi_{\mathrm{cyc}}|_\Delta\).
For a \(c\)-stable set \(\mathcal O\) of characters, put
\[
 e_{\mathcal O}=\sum_{\chi\in\mathcal O}e_\chi,\qquad
 \mathcal O^\dagger=\{\chi^\dagger:\chi\in\mathcal O\}.
\]
The dagger operation commutes with \(c\).
In the pair notation
\(\chi_{r,s}(\delta_1,\delta_2)=\delta_1^r\delta_2^s\),
where \(r-s\equiv0\pmod w\), these operations are exactly
\[
                  (r,s)^c=(s,r),\qquad
                  (r,s)^\dagger=(1-r,1-s)
                                \pmod{p-1}.
\]

\begin{proposition}\label{prop:cm-component-unit-detection}
Let \(\mathcal O\) be a \(c\)-stable set of characters of
\(\Delta\), and suppose that
\((e_{\mathcal O}\mathcal X_2)^+\) is finite.
Then
\[
 e_{\mathcal O}\mathcal J_2=0,\qquad
 e_{\mathcal O^\dagger}N_{\mathrm u}=0.
\]
In particular, the canonical abelian field
\[
 L_{\mathcal O}
   =\Omega^{\,\ker(H\longrightarrow A
                       \xrightarrow{\,e_{\mathcal O^\dagger}\,}
                                e_{\mathcal O^\dagger}A)}
\]
satisfies \(L_{\mathcal O}\subseteq L_{\mathrm u}\).
\end{proposition}

\begin{proof}
The idempotent \(e_{\mathcal O}\) commutes with \(c\), with
extension of ideals, and with ideal norm.
The norm maps on
\((e_{\mathcal O}\operatorname{Cl}(L_n)[p^\infty])^+\)
are consequently surjective.
The eventual-isomorphism argument in the preceding proof,
using \(N_{m/n}j_{m/n}=p^{2(m-n)}\), gives
\[
                      (e_{\mathcal O}\mathcal J_2)^+=0.
\]

We record the component map in the Kummer pairing, including
its twist.  For \(a\in N_{\mathrm u}\),
\(b\in\mathcal J_2\), and \(\delta\in\Delta\), equivariance is
\[
              \langle\delta a,\delta b\rangle
                    =\delta\langle a,b\rangle
                    =\omega_{\mathrm{cyc}}(\delta)
                                         \langle a,b\rangle,
\]
where the pairing is written additively.
Thus a \(\psi\)-component of \(N_{\mathrm u}\) pairs with a
\(\chi\)-component of \(\mathcal J_2\) exactly when
\(\psi\chi=\omega_{\mathrm{cyc}}\).
The isomorphism \eqref{eq:cm-unit-annihilator} therefore
restricts to the exact isomorphism
\[
 e_{\mathcal O^\dagger}N_{\mathrm u}
       \cong
 \operatorname{Hom}(e_{\mathcal O}\mathcal J_2,
                                      \mu_{p^\infty}).
\]
Because the source on the right has only the \(c\)-minus part,
the group on the left is \(c\)-fixed.
It is stable under \(\Gamma\), since \(\Delta\) commutes with
\(\Gamma\).
For \(a\) in this subgroup and \(\gamma\in\Gamma\), one has
\(\gamma a=c(\gamma a)=\gamma^c a\).
The elements \(\gamma(\gamma^c)^{-1}\) generate
\(\Gamma_{\mathrm a}\), so this subgroup lies in
\(A^{\Gamma_{\mathrm a}}\).
Lemma~\ref{lem:cm-anticyclotomic-invariants} makes it zero.
Perfect duality then gives \(e_{\mathcal O}\mathcal J_2=0\).

Finally the restriction kernel
\(\ker(H\to\operatorname{Gal}(L_{\mathrm u}/L_\infty))\)
maps into \(N_{\mathrm u}\) in \(A\).
Its \(e_{\mathcal O^\dagger}\)-projection is zero, so it is
contained in the kernel defining \(L_{\mathcal O}\).
The Galois correspondence proves
\(L_{\mathcal O}\subseteq L_{\mathrm u}\).
\end{proof}
